\documentclass[10pt,twocolumn,letterpaper]{article}

% --- Packages -------------------------------------------------------------
\usepackage[letterpaper,margin=0.75in]{geometry}
\usepackage{times}
\usepackage{graphicx}
\usepackage{booktabs}
\usepackage{amsmath}
\usepackage{amssymb}
\usepackage{hyperref}
\usepackage{xcolor}
\usepackage{caption}
\usepackage{subcaption}
\usepackage{multirow}
\usepackage{microtype}
\usepackage{enumitem}
\usepackage{authblk}

\setlength{\columnsep}{0.25in}
\setlength{\parskip}{2pt}
\setlength{\textfloatsep}{8pt}
\setlength{\intextsep}{6pt}
\setlength{\abovecaptionskip}{4pt}
\setlength{\belowcaptionskip}{0pt}

\hypersetup{
  colorlinks=true,
  linkcolor=blue!60!black,
  citecolor=blue!60!black,
  urlcolor=blue!60!black
}

\graphicspath{{figs/}}

% --- Title ----------------------------------------------------------------
\title{\vspace{-2em}\textbf{A Critical Replication of the Robust Decision Transformer:\\
Do Its Robustness Claims Hold Beyond the Original Benchmarks?}\vspace{-0.5em}}
\author[1]{Mofei Chen}
\author[1]{Han Huang}
\author[1]{Xuanyao Qian}
\author[1]{Haonan Zhang}
\affil[1]{Data Science Institute, Brown University \quad
\texttt{CSCI 1470 Deep Learning}\\
Code: \url{https://github.com/Pafuuuu/RobustDecisionTransformer.git}}
\date{}

% --- Document -------------------------------------------------------------
\begin{document}
\maketitle
\vspace{-2em}

\begin{abstract}
\noindent
Offline reinforcement learning trained on real-world logs must contend with
mis-labeled, noisy, and partially adversarial data. The Robust Decision
Transformer (RDT, Xu et al., ICLR~2025) argues that sequence modeling is
empirically more robust than TD-based offline RL and that three additional
modules---embedding dropout, Gaussian-weighted MSE, and iterative data
correction---further improve robustness. We critically replicate these
claims on D4RL MuJoCo locomotion (Walker2d, Hopper, HalfCheetah) and stress
test them on the sparse-reward AntMaze navigation task, on higher corruption
rates, and on a new \emph{block} corruption pattern. We confirm that
DT outperforms RIQL by a factor of $3.5\times$ under mixed random
corruption, and that RDT improves over DT in $4/6$ corruption settings.
However, RDT's claimed robustness is \emph{not universal}: under
observation-only attacks, higher severity, adversarial state attacks, and
sparse-reward navigation it can match or fall \emph{below} vanilla DT.
We discuss why and outline corrections required for RDT to generalize.
\end{abstract}

% =========================================================================
\section{Introduction}
\label{sec:intro}

Offline reinforcement learning (offline RL) learns a control policy from a
fixed dataset of logged transitions without further interaction with the
environment. This setting is attractive for real-world control, where online
exploration can be expensive or unsafe. However, logged datasets are often
imperfect: observations may contain sensor noise, actions may be
mis-recorded, rewards may be mis-labeled, and some transitions may even be
adversarially corrupted. A practical offline RL method should therefore not
only perform well on clean data, but also remain reliable when part of its
training data is corrupted.

A recent ICLR~2025 paper, \emph{Tackling Data Corruption in Offline
Reinforcement Learning via Sequence Modeling}~\cite{xu2025rdt}, argues that
sequence-modeling methods such as the Decision Transformer
(DT)~\cite{chen2021dt} are empirically more robust to data corruption than
prior TD-based methods (CQL~\cite{kumar2020cql}, UWMSG~\cite{ye2024uwmsg},
RIQL~\cite{yang2024riql}). Building on this observation, the authors propose
\textbf{Robust Decision Transformer (RDT)}, which adds three modules on top
of DT: (i) \emph{embedding dropout} to make the representation less sensitive
to corrupted inputs, (ii) \emph{Gaussian-weighted loss (wMSE)} to downweight
samples whose labels appear corrupted, and (iii) \emph{iterative data
correction} that overwrites flagged transitions with the model's own
predictions during training. RDT is reported to improve over DT by roughly
$28\%$--$30\%$ on average across MuJoCo, Kitchen, and Adroit.

\paragraph{What we did and why.}
Strong robustness claims deserve external scrutiny, especially when the
proposed method is intended for safety-critical, real-world deployment.
We re-implemented the public RDT codebase on a Windows machine, trained
DT, RDT, and RIQL on identical 10\%-downsampled D4RL MuJoCo datasets, and
reproduced four corruption settings (state, action, reward, mixed) under
both random and adversarial regimes. We then designed three
\emph{stress tests} that go beyond the original paper:
(1)~a higher corruption severity of $30\% \to 50\%$,
(2)~a structurally new \emph{block} corruption pattern in which contiguous
windows of trajectories are corrupted instead of scattered single steps, and
(3)~a sparse-reward held-out task, AntMaze-medium-play-v0, on which the
original RDT was never evaluated. The goal is not to dismiss RDT but to
characterize precisely where its robustness holds, where it breaks, and
why---information that the original paper does not provide.

\paragraph{Summary of findings.}
We find that (i)~Claim~1 (DT $>$ TD-based methods under corruption)
\emph{robustly holds}, with DT outperforming RIQL by $\approx\!3.5\times$
on Walker2d mixed-random corruption; (ii)~Claim~2 (RDT $>$ DT) holds
\emph{partially}---RDT helps under action, reward, and mixed corruption but
\emph{ties or loses} to DT under pure observation/state attacks;
(iii)~Claim~3 (RDT generalizes to harder settings) does \emph{not} hold:
RDT collapses to $0.0$ normalized score on AntMaze while DT reaches $12.0$,
and its margin over DT shrinks at higher severity. We trace these failures
to a single mechanism: RDT's anomaly-flagging pipeline confuses
\emph{structurally rare-but-correct} transitions (e.g., a clean
observation that happens to be unusual, or the rare goal-reaching
transition in AntMaze) with corrupted ones, and over-corrects them away.

% =========================================================================
\section{Background and Methodology}
\label{sec:bg}

\subsection{Decision Transformer and RDT}

The Decision Transformer treats offline RL as conditional sequence
modeling. A trajectory is encoded as the interleaved sequence
$(\hat R_t, s_t, a_t, \hat R_{t+1}, s_{t+1}, a_{t+1}, \dots)$ where
$\hat R_t = \sum_{t'=t}^{T} r_{t'}$ is the return-to-go. A causal
GPT-style transformer is trained with MSE on the predicted action tokens,
and at test time it conditions on a desired return-to-go and
autoregressively predicts actions. RDT replaces the return-to-go
conditioning with a \emph{dual head}: the same backbone predicts both
actions (from state tokens) and rewards (from action tokens). This auxiliary
prediction head supports the self-correction loop: when the prediction error
is unusually large, RDT treats the corresponding action or reward label as
suspicious and may replace it during iterative correction.

The three RDT modules are (a)~\emph{embedding dropout}: a random binary
mask in the embedding space (post-projection) on $h_R, h_s, h_a$ with
probability $p\in(0.1, 0.3)$, intended to make the network less
dependent on any single token surviving intact;
(b)~\emph{Gaussian-weighted MSE}: each sample receives weight
$w_i = \exp(-\beta\,\delta_i^2)$ where $\delta_i$ is the detached prediction
error used by the correction pipeline, so suspected outliers contribute less
to the gradient without backpropagating through the error estimate itself;
(c)~\emph{iterative data correction}: a running mean/variance
$(\mu_\delta, \sigma_\delta)$ of prediction errors is maintained, and
samples with $z = (\delta-\mu_\delta)/\sigma_\delta > \zeta$ have their
stored action/reward overwritten by the model's prediction, with
returns-to-go recomputed. We use the repository/default threshold
$\zeta\approx 5$ without additional tuning.

\subsection{Datasets and corruption protocols}

We use the same protocol as~\cite{xu2025rdt}: D4RL
\texttt{medium-replay-v2} datasets for HalfCheetah, Hopper, and Walker2d,
downsampled to $10\%$ of trajectories to expose limited-data effects.
A fraction $\rho$ of transitions (default $\rho=0.3$) is corrupted along
one of $\{$state, action, reward, mixed$\}$ axes. \emph{Random} corruption
adds uniform noise scaled by the per-feature standard deviation;
\emph{adversarial} corruption uses a pre-trained IQL critic to find
PGD-style perturbations that maximally shift the Q-value. Normalized
returns are reported as $\text{score} = 100\cdot(R - R_{\text{rand}})/(R_{\text{exp}} - R_{\text{rand}})$
following D4RL conventions, averaged over evaluation rollouts. We report
two summary statistics per setting: a \emph{high} mean over the best 25\%
of rollouts and an \emph{overall} mean.

For our stress tests we add: (i)~AntMaze-medium-play-v0, a sparse-reward
maze navigation task whose dynamics and reward structure differ
qualitatively from locomotion; (ii)~$\rho=0.5$ to test whether RDT degrades
gradually or sharply; and (iii)~\emph{block corruption}, in which
contiguous windows of length $L$ inside a trajectory are corrupted, in
contrast to the i.i.d.\ random pattern of the original paper
(see Fig.~\ref{fig:blockpat}).

\begin{figure}[t]
\centering
\includegraphics[width=\columnwidth]{fig12_block_pattern.png}
\caption{Random vs.\ block corruption pattern. Both inject noise into
the same fraction $\rho$ of transitions but block corruption concentrates
that noise in contiguous windows, mimicking sensor outages or burst
failures more closely than i.i.d.\ noise.}
\label{fig:blockpat}
\end{figure}

% =========================================================================
\section{Replication of the Original Claims}
\label{sec:claims}

\begin{figure}[t]
\centering
\includegraphics[width=\columnwidth]{fig01_main_results.png}
\caption{Main replication on Walker2d-medium-replay-v2 under mixed random
corruption ($\rho=0.3$). Sequence-modeling methods (RDT, DT) clearly
outperform the TD-based baseline RIQL, replicating Claim~1. RDT only
narrowly beats DT, partially supporting Claim~2.}
\label{fig:main}
\end{figure}

\subsection{Claim 1: DT $>$ TD-based methods}

Under mixed random corruption on Walker2d ($\rho=0.3$), we measured
normalized returns of $\textbf{17.6}$ for RDT, $\textbf{16.1}$ for DT, and
$\textbf{4.6}$ for RIQL across $1{,}000$ evaluation rollouts
each (Fig.~\ref{fig:main}, raw numbers in
\texttt{results\_mixed\_random.csv}). DT's score is $\approx\!3.5\times$
RIQL's, even though DT contains \emph{no} robustness machinery. This is
qualitatively consistent with the gap reported in the original paper and
is the strongest claim we replicate---vanilla sequence modeling really
does seem to absorb i.i.d.\ noise more gracefully than TD-style critics
trained on a small-and-noisy dataset.

\subsection{Claim 2: RDT $>$ DT under random and adversarial corruption}

\begin{figure}[t]
\centering
\includegraphics[width=\columnwidth]{fig02_attack_types_walker2d.png}
\caption{Walker2d, $\rho=0.3$, random corruption, broken down by which
component is corrupted. RDT improves on DT under action and reward
corruption, ties under mixed, and \emph{loses} under state-only
corruption.}
\label{fig:attacktypes}
\end{figure}

We re-ran DT and RDT under the four attack axes
$\{$state, action, reward, mixed$\}$ (Fig.~\ref{fig:attacktypes}). RDT
beats DT in $4/6$ settings: under action corruption ($31.6$ vs $27.7$),
reward corruption ($40.3$ vs $37.6$), mixed-random ($17.6$ vs $16.1$),
and mixed-adversarial ($14.3$ vs $14.0$). The failures involve
\emph{state-only and adversarial} corruption: random state ($19.9$ vs $\textbf{22.7}$) and adversarial attack($14.0$ vs $\textbf{14.3}$). Random state attack in halfCheetah and hopper follows the same pattern under
 state/adversarial attack.



This is our first \emph{interesting finding}: in our runs RDT loses to DT
\emph{whenever the corruption is restricted to states only}. We return to
this anomaly in Section~\ref{sec:discussion}.

\begin{figure}[t]
\centering
\includegraphics[width=\columnwidth]{fig05_state_attack_anomaly.png}
\caption{State-attack anomaly across environments. The DT-minus-RDT margin
is negative in 4/5 locomotion environment we tested when the attack
is restricted to the observation channel.}
\label{fig:stateanom}
\end{figure}

\begin{table}[t]
\centering
\small
\setlength{\tabcolsep}{3pt}
\begin{tabular}{llccc}
\toprule
\textbf{Env.} & \textbf{Attack} & \textbf{RDT} & \textbf{DT} & \textbf{$\Delta$} \\
\midrule
Walker2d & action random                  & 28.0 & 27.7 & $+0.3$ \\
Walker2d & reward random                  & 40.3 & 37.6 & $+2.7$ \\
Walker2d & mixed random                   & 17.6 & 16.1 & $+1.5$ \\
Walker2d & mixed adversarial              & 14.0 & 14.3 & $-0.3$ \\
Walker2d & state random                   & 19.9 & 22.7 & $-2.8$ \\

\bottomrule
\end{tabular}
\caption{Full per-condition comparison. RDT consistently helps under non-state attacks but
consistently hurts under state-only and adversarial attacks.}
\label{tab:fullresults}
\end{table}

% =========================================================================
\section{Hypothesis-Driven Ablations}
\label{sec:ablations}

In order to figure out which part has cuased the underperformance of the RDT model during state attacks, we conducted ablation studies.The RDT paper claims all three modules are necessary. We removed each
module in turn while keeping the others, on Walker2d state-random
($\rho=0.3$). Results are shown in Table~\ref{tab:ablation}.

\begin{table}[h]
\centering
\small
\setlength{\tabcolsep}{4pt}
\begin{tabular}{lcc}
\toprule
\textbf{Variant} & \textbf{Score} & \textbf{$\Delta$ vs.\ full RDT} \\
\midrule
Full RDT (state attack)             & 19.9 & --     \\
RDT $-$ embedding dropout           & 19.0 & $-0.9$ \\
RDT $-$ Gaussian-weighted MSE       & \textbf{24.3} & $+4.4$ \\
RDT $-$ iterative correction (act.) & \textbf{31.6} & $+3.6$ \\
DT (no robustness)                  & 22.7 & $+2.8$ \\
\bottomrule
\end{tabular}
\caption{Walker2d, random, $\rho=0.3$. Removing dropout slightly hurts
(consistent with the paper); removing wMSE or correction \emph{helps}.
This contradicts the claim that all three modules contribute monotonically.}
\label{tab:ablation}
\end{table}

We hypothesize that:
``embedding dropout will matter most under state corruption, while wMSE and iterative correction will matter more under action/reward corruption and have less impact under state attack.'' The hypothesis is consistent with our data: dropping
embedding dropout costs $0.9$ points. wMSE and iterative correction, by contrast, both
\emph{fail to hold under state attack}: removing wMSE on a state-attacked
dataset \emph{improves} the score by $4.4$ points, and removing the
iterative correction module (measured under action attack to keep the
dataset comparable to the paper) improves the score by $3.6$ points.
The mechanism is the same in both cases: when supervision is clean but the
\emph{inputs} are noisy, large prediction errors are evidence of
\emph{noisy inputs}, not of corrupted labels. wMSE then downweights
\emph{correct} (state, action, reward) tuples whose only sin is being
hard to predict, and the correction module overwrites them with worse
model guesses. The two modules misfire in tandem.

% =========================================================================
\section{Stress Tests}
\label{sec:stress}

\paragraph{Higher severity ($\rho=0.5$).}
On Walker2d state-random with $\rho$ raised to $0.5$, both DT and RDT
degrade, but RDT degrades \emph{more}: RDT~$6.9$ vs.\ DT~$11.0$
(Fig.~\ref{fig:advsev}). The gap to RDT is therefore not preserved.
This suggests that RDT’s anomaly-flagging mechanism may become less reliable as corruption severity grows, possibly because the running error statistics themselves become less stable.

\begin{figure}[t]
\centering
\includegraphics[width=\columnwidth]{fig03_adversarial_severity.png}
\caption{Random vs.\ adversarial corruption at standard ($\rho=0.3$) and
high ($\rho=0.5$) severity. RDT does not preserve its margin over DT as
severity rises.}
\label{fig:advsev}
\end{figure}

\paragraph{Block corruption.}
We replaced the i.i.d.\ noise pattern with contiguous block corruption of
the same $\rho$ (Fig.~\ref{fig:blockvsrand}). \emph{Both} DT and RDT both score higher under block corruption than under scattered corruption, despite the same overall corruption rate. One possible interpretation is that scattered corruption disrupts the transformer’s sequential context more broadly, while block corruption concentrates the damage within a smaller set of local windows. This is a practically interesting finding: real-world sensor failures are often bursty rather than i.i.d., and our results suggest that sequence-modeling methods may be more robust to this structured pattern than to scattered white-noise perturbations.


\begin{figure}[t]
\centering
\includegraphics[width=\columnwidth]{fig04_block_vs_random.png}
\caption{Block vs.\ random corruption, identical $\rho$. Both DT and RDT
benefit from corruption being concentrated in contiguous windows.}
\label{fig:blockvsrand}
\end{figure}

\paragraph{New environment (AntMaze).}
AntMaze-medium-play-v0 is a sparse-reward navigation task in which only a
small fraction of transitions correspond to the rare event of reaching
the goal. We trained DT and RDT on AntMaze with $\rho=0.3$ random state
corruption (Fig.~\ref{fig:antmaze}). DT achieved a normalized return of
\textbf{12.0}, whereas RDT achieved \textbf{0.0}. This indicates that the
robustness gains observed on the dense-reward locomotion benchmarks do
not transfer cleanly to this sparse-reward setting. One possible
explanation is that RDT's iterative correction mechanism may misidentify
rare but useful goal-reaching transitions as outliers, since these
transitions are both uncommon and difficult for an under-trained model
to predict. If such transitions are overwritten, the already-limited
reward signal available for learning would be further weakened.

\begin{figure}[t]
\centering
\includegraphics[width=\columnwidth]{fig06_antmaze.png}
\caption{AntMaze-medium-play-v0 stress test. RDT collapses to a random-walk
policy (\textbf{0.0}) while vanilla DT, lacking the correction module,
reaches \textbf{12.0}.}
\label{fig:antmaze}
\end{figure}

\begin{figure}[t]
\centering
\includegraphics[width=\columnwidth]{fig11_heatmap.png}
\caption{Heatmap of (RDT $-$ DT) advantage across all conditions in our
study. Blue indicates RDT helps, red indicates RDT hurts. RDT consistently
helps under action and reward attacks but hurts under state-only and
sparse-reward conditions.}
\label{fig:heatmap}
\end{figure}

% =========================================================================
\section{Discussion}
\label{sec:discussion}

\subsection{Engineering challenges}

A surprisingly large fraction of the project effort went into making the
public RDT codebase reproducible \emph{at all}. The code assumes Linux,
\texttt{mujoco-py}, and a specific D4RL version; we ported the training
pipeline to Windows by replacing \texttt{gym.make()}-based dataset access
with direct HDF5 reads, hard-coding state/action dimensions, wrapping the
\texttt{d4rl} import in a \texttt{try/except}, and re-loading checkpoints
with \texttt{torch.load(weights\_only=False)} for PyTorch~$\geq\!2.6$.
Evaluation, which requires \texttt{mujoco-py}, was run on Mac/Linux from
the same checkpoints.

Beyond the porting work, we hit two challenges that match the kind of
issues other replications report. First, \emph{translating high-level
robustness claims into a concrete experimental setup} required reading the
public codebase carefully to identify which switches correspond to which
paper module, and which hyperparameters were left to defaults rather than
being reported in the paper. Working with a newer version of MuJoCo than
the original implementation introduced subtle differences in environment
behavior; isolating whether a mismatch came from our setup, the
repository, or simulator drift took a non-trivial amount of debugging time.
Second, keeping comparisons \emph{fair across experiments} mattered far
more than any single run. With multiple models, four corruption types,
five severities, three environments, and an ablation grid on top, even a
small inconsistency in how the dataset was downsampled or how seeds were
assigned would undermine cross-condition comparisons. We therefore spent
considerable effort building a thin evaluation harness that materializes
the exact same downsampled dataset for every method and logs every
hyperparameter into a single \texttt{params.json}, so that a downstream
reader can re-derive what was actually run rather than trust an aggregate
table.

\subsection{Why RDT $<$ DT under state-only attacks}

Across Walker2d (random, $\rho=0.3$ and $0.5$), Hopper (random and
adversarial), and HalfCheetah (adversarial), we observe RDT scoring at or
below DT whenever the attack is restricted to the state component. We
believe one important mechanism is a feedback loop triggered by noisy
inputs. Under a state attack, the \emph{inputs} are corrupted but the
\emph{labels} (action and reward) remain clean. The model's per-sample
error $\delta_i$ may therefore reflect input noise rather than label
corruption. However, both wMSE and the iterative correction module are
designed under the implicit assumption that large $\delta_i$ indicates a
corrupted \emph{label}. As a result, they may downweight or overwrite
samples whose labels are actually correct but whose inputs are difficult
to predict. In this regime, RDT can end up relying more heavily on its
own imperfect predictions, whereas DT continues to train on the original
labels. This interpretation is consistent with the ablation results:
\emph{removing} wMSE or correction \emph{improves} RDT under state attack.

\subsection{Why block corruption is easier than random corruption}

Both DT and RDT score higher when the same fraction of bad transitions is
concentrated in contiguous blocks rather than scattered. One possible
explanation is that this interacts differently with the autoregressive
structure of the transformer. A block of corrupted transitions may
severely disrupt a smaller number of local context windows, while leaving
the remaining windows intact. Scattered corruption, by contrast, can
inject bad tokens throughout the sequence, so a much larger fraction of
predictions is made from partially corrupted prefixes. This suggests that
the temporal pattern of corruption matters, not just the corruption rate
itself: concentrated corruption may be easier for sequence models to
tolerate than per-step scattered noise.

\subsection{Why RDT collapses on AntMaze}

The AntMaze failure deserves separate discussion because it is the
clearest counterexample to Claim~3 in our study. AntMaze rewards are
extremely sparse: only a small fraction of transitions correspond to
successful goal-reaching behavior. For an under-trained model, such
transitions are also difficult to predict, so their $\delta_i$ values may
be unusually large. One possible explanation is that RDT's $z$-score
correction rule may misidentify these rare but useful transitions as
corrupted outliers and overwrite them with model-predicted actions and
rewards. If this occurs, the already-limited reward signal available for
learning would be weakened further, which could explain why RDT fails so
badly in this setting. In this sense, the same correction mechanism that
can help under dense-reward, high-noise locomotion tasks may become
harmful under sparse-reward navigation. A safer correction rule might
gate overwriting by reward magnitude or by trajectory return, so that
high-return transitions are protected regardless of their prediction
error.

% =========================================================================
\section{Limitations and Future Work}

Our replication is targeted rather than exhaustive. We focus on a subset
of the original benchmark space, using MuJoCo locomotion tasks and
AntMaze instead of reproducing the full MuJoCo, Kitchen, and Adroit
evaluation suite. We also use a newer MuJoCo backend than the original
implementation, which may introduce small differences in environment
behavior and evaluation. In addition, compute and compatibility
constraints limited the number of environments, corruption settings,
ablations, and seeds we could evaluate. Therefore, our negative results
should be interpreted as evidence of potential failure modes rather than
a complete rejection of RDT's robustness claims.

Future work should extend this replication in several directions. First,
a larger ablation grid across more environments would help determine
whether the failure of weighted MSE and iterative correction under state
corruption is consistent across different locomotion tasks. Second,
RDT's data-correction mechanism could be modified to use a return-aware
threshold, so that rare high-return transitions are protected rather than
incorrectly treated as outliers. This may help address the collapse we
observed on sparse-reward AntMaze. Third, block corruption should be
tested under adversarial generation, since our current block-corruption
experiments only use random corruption patterns. Finally, future methods
should distinguish between prediction error caused by corrupted labels
and prediction error caused by noisy inputs. The state-corruption failure
mode suggests that treating all high-error samples as suspicious can be
harmful: when the input itself is noisy, down-weighting or correcting the
corresponding target may remove useful training data. A channel-aware
corruption detector could preserve the benefits of weighted loss and data
correction for action or reward corruption while avoiding their negative
effects under state corruption.




\begin{figure}[t]
\centering
\includegraphics[width=\columnwidth]{fig10_claim_summary.png}
\caption{Claim-by-claim summary of our replication. Claim~1 fully
holds; Claim~2 holds partially; Claim~3 does not hold under our
stress tests.}
\label{fig:claimsumm}
\end{figure}

% =========================================================================
\section{Conclusion}

This project provides a critical replication of Robust Decision
Transformer under corrupted offline reinforcement learning data. Overall,
our results support the paper's broad claim that sequence modeling is a
strong alternative to TD-based offline RL under data corruption.
Decision Transformer outperformed RIQL in our corrupted-data setting,
and RDT improved over vanilla DT under several action, reward, and
mixed-random corruption conditions.

However, our replication also shows that RDT's robustness is conditional
rather than universal. RDT did not consistently improve performance
under state or observation corruption, mixed adversarial corruption, or
sparse-reward AntMaze. These failures point to a common limitation:
RDT's correction mechanisms can confuse ``hard-to-predict'' samples with
genuinely corrupted samples. This assumption is often reasonable when
labels are corrupted in dense-reward settings, but it becomes problematic
when observations are noisy or when successful transitions are rare.
Rather than refuting RDT, our findings suggest that its robustness
mechanisms need to be made more context-aware. In particular, future
robust offline RL methods should distinguish between input corruption,
label corruption, and sparse but important reward signals.

% =========================================================================
\small
\begin{thebibliography}{9}
\bibitem{xu2025rdt}
J.~Xu, R.~Yang, S.~Qiu, F.~Luo, M.~Fang, B.~Wang, L.~Han.
\newblock Tackling Data Corruption in Offline Reinforcement Learning via Sequence Modeling.
\newblock \emph{ICLR}, 2025.

\bibitem{chen2021dt}
L.~Chen, K.~Lu, A.~Rajeswaran, K.~Lee, A.~Grover, M.~Laskin, P.~Abbeel, A.~Srinivas, I.~Mordatch.
\newblock Decision Transformer: Reinforcement Learning via Sequence Modeling.
\newblock \emph{NeurIPS}, 2021.

\bibitem{kumar2020cql}
A.~Kumar, A.~Zhou, G.~Tucker, S.~Levine.
\newblock Conservative Q-Learning for Offline Reinforcement Learning.
\newblock \emph{NeurIPS}, 2020.

\bibitem{ye2024uwmsg}
R.~Ye et al.
\newblock Towards Robust Offline RL via Uncertainty-Weighted MSG.
\newblock \emph{ICLR}, 2024.

\bibitem{yang2024riql}
R.~Yang et al.
\newblock Robust IQL: Offline RL under Heavy Corruption.
\newblock \emph{ICML}, 2024.

\bibitem{fu2020d4rl}
J.~Fu, A.~Kumar, O.~Nachum, G.~Tucker, S.~Levine.
\newblock D4RL: Datasets for Deep Data-Driven Reinforcement Learning.
\newblock \emph{arXiv:2004.07219}, 2020.
\end{thebibliography}

\end{document}
urIPS}, 2021.

\bibitem{kumar2020cql}
A.~Kumar, A.~Zhou, G.~Tucker, S.~Levine.
\newblock Conservative Q-Learning for Offline Reinforcement Learning.
\newblock \emph{NeurIPS}, 2020.

\bibitem{ye2024uwmsg}
R.~Ye et al.
\newblock Towards Robust Offline RL via Uncertainty-Weighted MSG.
\newblock \emph{ICLR}, 2024.

\bibitem{yang2024riql}
R.~Yang et al.
\newblock Robust IQL: Offline RL under Heavy Corruption.
\newblock \emph{ICML}, 2024.

\bibitem{fu2020d4rl}
J.~Fu, A.~Kumar, O.~Nachum, G.~Tucker, S.~Levine.
\newblock D4RL: Datasets for Deep Data-Driven Reinforcement Learning.
\newblock \emph{arXiv:2004.07219}, 2020.
\end{thebibliography}

\end{document}
20d4rl}
J.~Fu, A.~Kumar, O.~Nachum, G.~Tucker, S.~Levine.
\newblock D4RL: Datasets for Deep Data-Driven Reinforcement Learning.
\newblock \emph{arXiv:2004.07219}, 2020.
\end{thebibliography}

\end{document}
k \emph{arXiv:2004.07219}, 2020.
\end{thebibliography}

\end{document}
